% ===================================================================
%  সারণি নং 15 — বাজার। — খাদ্যসামগ্রী।
%
%  Ceci n'est PAS une transcription : c'est une traduction, faite en
%  2026, en bengali d'aujourd'hui. Voir texto/bn/10-sarani-01.tex
%  pour la consigne, qui vaut pour les seize fichiers.
% ===================================================================

%%K t15-apar-1 apar
\VUsaut{4.0mm}
\VUtitre{15.0pt}{60}{সারণি নং 15}
\VUsaut{3.0mm}

%%K t15-tit-1 sub
\VUpk{11.4pt}{40}{বাজার। --- খাদ্যসামগ্রী।}
\VUsaut{3.0mm}

%%K t15-01-1 p
1. --- যে ব্যবসায়ীরা আমাদের প্রয়োজনের খাবার জোগান, তাঁদের অধিকাংশ
\VUgras{ঢাকা বাজারের} \VUgras{হলের} \textsuperscript{(1)} নিচে জড়ো থাকেন, বা
পাশের গলিতে।

%%K t15-02-1 p
2. --- \VUgras{শুয়োরের মাংসের কসাই} \textsuperscript{(2)}, যে \VUgras{হ্যাম}
\textsuperscript{(3)}, \VUgras{রক্তের সসেজ} \textsuperscript{(4)},
\VUgras{সসেজ} \textsuperscript{(5)}, \VUgras{অন্ত্রের সসেজ} \textsuperscript{(6)}
ও \VUgras{চর্বি} \textsuperscript{(7)} বেচে, সে \VUgras{মাখন ও পনির
বিক্রেতার} \textsuperscript{(8)} পাশে বসে, যার ঠেলা লাল খোসাওয়ালা
\VUgras{হল্যান্ডের পনিরে} \textsuperscript{(9)}, \VUgras{কামঁবেরে}
\textsuperscript{(10)}, \VUgras{গ্রুইয়েরের বড় চাকতিতে} \textsuperscript{(11)} ও
টাটকা \VUgras{মাখনের তালে} \textsuperscript{(12)} সাজানো।

%%K t15-03-1 p
3. --- তার সামনে এক \VUgras{তরকারি-বিক্রেতা} \textsuperscript{(13)} নিজের
খদ্দেরদের সুন্দর \VUgras{টমেটো} \textsuperscript{(14)}, \VUgras{ফুলকপি}
\textsuperscript{(15)}, \VUgras{লেটুস} \textsuperscript{(16)}, \VUgras{বাঁধাকপি}
\textsuperscript{(17)}, \VUgras{কুমড়ো} \textsuperscript{(19)}, \VUgras{খরমুজ}
\textsuperscript{(18)} এমনকি \VUgras{ছত্রাক} \textsuperscript{(20)} পর্যন্ত
দেখায়। কিন্তু এক \VUgras{মদওয়ালা} নিজের \VUgras{টানার গাড়ি}
\textsuperscript{(21)}, যা \VUgras{বিয়ারের পিপেয়} \textsuperscript{(22)} লদপদ,
ঠিক তার ঠেলার সামনে দাঁড় করিয়ে দিয়েছে আর তার খদ্দেরদের কাছে আসতে দিচ্ছে না।
এক মালি তাকে \VUgras{আর্টিচোকের} \textsuperscript{(23)} একটি ঝুড়ি আনছে।

%%K t15-tit-2 sub
\VUsaut{4.0mm}
\VUpk{10.2pt}{40}{বেচা ও কেনা।}
\VUsaut{3.0mm}

%%K t15-04-1 p
4. --- বাজারে বড় সাড়া, কারণ \VUgras{নিলামের} \textsuperscript{(24)} সময় এসে
গেছে। যে \VUgras{গৃহিণীরা} \textsuperscript{(26)} এখানে বাজার করতে আসেন, তাঁরা
পথে \VUgras{ভুঁড়ি-বিক্রেতার} \textsuperscript{(25)} ঠেলার সামনে থেমে
\VUgras{ভুঁড়ি} বা একটি \VUgras{বাছুরের মাথা} কেনেন।

%%K t15-05-1 p
5. --- এক মোটা \VUgras{রাঁধুনি} \textsuperscript{(27)} \VUgras{মাছ-বিক্রেতার}
\textsuperscript{(28)} সঙ্গে একটি খুব সুন্দর \VUgras{টুনার} দর কষছে, যিনি নিজের
খুশিমতো দাম বাড়িয়ে চলেছেন। তাঁর কাছে \VUgras{বান, সোল, ট্রাউট} ও
\VUgras{কার্পের} বড় চালান এসেছে। তাঁর \VUgras{কড} ও তাঁর \VUgras{ঝিনুক} খুব
টাটকা।

%%K t15-tit-3 sub
\VUsaut{4.0mm}
\VUpk{10.2pt}{40}{সস্তা মাল ও দামি।}
\VUsaut{3.0mm}

%%K t15-06-1 p
6. --- তাঁর পাশে বসা \VUgras{মুরগি-বিক্রেতা} \textsuperscript{(29)} খুব সৎ। তিনি
যখন কোনো \VUgras{টার্কি}, কোনো \VUgras{রাজহাঁস}, কোনো \VUgras{হাঁস} বা সাধারণ
\VUgras{মুরগি} বেচেন, তখন ন্যায্য দামের বেশি চান না।

%%K t15-07-1 p
7. --- চত্বরের কোণে, প্রথম তলায়, এক \VUgras{কাপড়ের দোকানদার}
\textsuperscript{(30)} সম্প্রতি দোকান খুলেছেন, যেখানে ক্রেতারা
\VUgras{টেবিলক্লথ}, \VUgras{ন্যাপকিন, এপ্রন} ও \VUgras{তোয়ালের} খুব চমৎকার
বাছাই সব সময় পান। সেই তলাতেই এক \VUgras{ছুরি-নির্মাতা} \textsuperscript{(31)}
নিজের কর্মশালা বসিয়েছেন। তিনি বাজারের মেয়েদের \VUgras{ছুরি} ধার দেন আর মালির
জন্য সবচেয়ে কড়া \VUgras{ইস্পাত} দিয়ে \VUgras{ছাঁটার কাস্তে} বানান।

%%K t15-tit-4 sub
\VUsaut{4.0mm}
\VUpk{10.2pt}{40}{মুদির দোকান।}
\VUsaut{3.0mm}

%%K t15-08-1 p
8. --- তাঁর কর্মশালার নিচে কিছুকাল আগে একটি বড় খাদ্যদোকান খুলেছে। এটি আর
পুরনো দিনের সেই সাদাসিধা ছোট \VUgras{মুদির দোকান} \textsuperscript{(32)} নয়, যা
কেবল রোজকার মাল বেচত --- \VUgras{চা} \textsuperscript{(33)}, \VUgras{চকোলেট}
\textsuperscript{(34)}, \VUgras{ডেলা চিনি} \textsuperscript{(37)} ও \VUgras{সাবান}
\textsuperscript{(38)}। এখানে সব কিছু বিকোয়: \VUgras{মুচমুচে বিস্কুট},
\VUgras{কৌটোবন্দি জিনিস} \textsuperscript{(35)}, \VUgras{মদ} \textsuperscript{(36)},
\VUgras{রাম, ব্র্যান্ডি, কির্শ, অ্যানিসেট} ও \VUgras{কুরাসাও}। এখানে শিকারও
ততটাই সহজে মেলে --- \VUgras{খরগোশ} \textsuperscript{(39)}, \VUgras{পাখি}
\textsuperscript{(40)}, \VUgras{তিতির} \textsuperscript{(41)}, \VUgras{বটের}
\textsuperscript{(42)}, \VUgras{বুনো হাঁস} \textsuperscript{(43)} বা
\VUgras{ময়ূর-তিতির} \textsuperscript{(44)} --- যতটা গরম দেশের ফল, যেমন
\VUgras{কলা} \textsuperscript{(46)} ও \VUgras{আনারস}।

%%K t15-09-1 p
9. --- এক খুব উপকারী মুদির \VUgras{কর্মচারী} \textsuperscript{(45)} যা চাওয়া হয়
তা-ই দিতে তৈরি থাকে: \VUgras{মোরব্বা} \textsuperscript{(47)}, বেদানার বা বেলের,
\VUgras{চর্বি} \textsuperscript{(48)}, \VUgras{শসা} \textsuperscript{(49)} বা
নোনতা \VUgras{সার্ডিন} \textsuperscript{(50)}।

%%K t15-09-2 p
এটি \VUgras{ক্রেত্রীদের} \textsuperscript{(51)} পক্ষে খুব সুবিধার কথা, যাঁরা
সবচেয়ে সাধারণ মালের পাশেই সবচেয়ে দুর্লভ আগাম ফসল ও সবচেয়ে সুস্বাদু ফল পান:
রসালো \VUgras{নাশপাতি} \textsuperscript{(52)}, \VUgras{আলুবোখারা}
\textsuperscript{(53)}, \VUgras{আঙুর} \textsuperscript{(54)}, \VUgras{পিচ}
\textsuperscript{(55)}, \VUgras{বাদাম} \textsuperscript{(56)} ও \VUgras{আখরোট}
\textsuperscript{(57)}।

%%K t15-tit-5 sub
\VUsaut{4.0mm}
\VUpk{10.2pt}{40}{খদ্দেরেরা।}
\VUsaut{3.0mm}

%%K t15-10-1 p
10. --- \VUgras{চাল} \textsuperscript{(58)} ও \VUgras{মসুর} \textsuperscript{(59)}
ধোঁয়ানো \VUgras{হেরিঙের} \textsuperscript{(60)} বাক্সের পাশে রাখা;
\VUgras{আলু} \textsuperscript{(61)} ও \VUgras{শিম} \textsuperscript{(63)}
\VUgras{আপেলের} \textsuperscript{(62)}, \VUgras{ডুমুরের} \textsuperscript{(64)},
\VUgras{শুকনো আলুবোখারার} \textsuperscript{(65)} ও \VUgras{পাহাড়ি বাদামের}
\textsuperscript{(66)} পাশে বসে। এই দোকানে টাটকা \VUgras{ডিম}
\textsuperscript{(67)} ও \VUgras{জলপাইও} \textsuperscript{(68)} মেলে, তাই
\VUgras{ঝুড়ি} \textsuperscript{(70)} ও \VUgras{জালের থলি} \textsuperscript{(73)}
খুব তাড়াতাড়ি ভরে যায়।

%%K t15-11-1 p
11. --- এই চমৎকার দোকানের \VUgras{কফি} \textsuperscript{(72)}
\VUgras{কাজের মেয়েদের} \textsuperscript{(69)} ও রাঁধুনিদের মধ্যে বেশ ভালো নাম
কামিয়েছে, কারণ বুড়ো \VUgras{মুদি} \textsuperscript{(71)} সেটি নিজেই বড়
যত্নে ভাজেন।

%%K t15-tit-6 sub
\VUsaut{4.0mm}
\VUpk{10.2pt}{40}{দেখার জিনিস।}
\VUsaut{3.0mm}

%%K t15-12-1 p
12. --- সবাই বাজারে বাজার করতেই আসে না। হল পর্যন্ত যাওয়া সেই ছোট গলির মনোরম
আনাগোনায় নানা রকম লোক চোখে পড়ে। এক ভালোমানুষ \VUgras{কসাইখানার মজুরের}
\textsuperscript{(75)} পাশে, যে নিজের আগে একটি \VUgras{ঠেলা} \textsuperscript{(74)}
ঠেলছে, এই যে ছুটিতে আসা দুই সিপাহি, যাদের নিজেদের ঝলমলে উর্দি দেখাতে বড় গর্ব:
\VUgras{হুসার} \textsuperscript{(76)} নিজের নীল \VUgras{দোলমানে} ও
\VUgras{পালকওয়ালা শাকোয়}, আর \VUgras{জুয়াভের হাবিলদার} ঢিলে নিকারে, মাথায়
\VUgras{ফেজ} পরে।

%%K t15-13-1 p
13. --- \VUgras{রুটিওয়ালা} \textsuperscript{(80)}, যে \VUgras{রুটির দোকানের}
\textsuperscript{(78)} দোরগোড়ায় দাঁড়িয়ে, সে এইমাত্র নিজের \VUgras{রুটির}
\textsuperscript{(79)} মণ্ড মেখেছে ও তন্দুর থেকে সেই \VUgras{ক্রোয়াসঁ}
\textsuperscript{(81)}, \VUgras{ছোট রুটি} \textsuperscript{(82)} ও \VUgras{গোল
রুটি} \textsuperscript{(83)} বের করেছে যা জানালায় রাখা।

%%K t15-14-1 p
14. --- রুটির দোকানের উপরে এক দক্ষ \VUgras{সেলাইকারিণী} \textsuperscript{(84)}
থাকেন, যিনি কয়েকজন \VUgras{সূচিশিল্পিনী} \textsuperscript{(85)} রাখেন। তাঁর
পাশে এক \VUgras{ধোপানি ও ইস্তিরিওয়ালি} \textsuperscript{(86)} থাকেন, যিনি
\VUgras{কাপড়} \textsuperscript{(88)} ধুয়ে ও শুকিয়ে তাতে মাড় দেন আর
\VUgras{ইস্তিরি} \textsuperscript{(87)} দিয়ে চাপেন।

%%K t15-tit-7 sub
\VUsaut{4.0mm}
\VUpk{10.2pt}{40}{মাংস, মাছ ও তরকারি।}
\VUsaut{3.0mm}

%%K t15-15-1 p
15. --- নিচের তলার \VUgras{কসাইয়ের দোকানে} \textsuperscript{(89)} সব রকম মাংস
বিকোয় --- \VUgras{ভেড়ার} \textsuperscript{(90)}, \VUgras{গরুর}
\textsuperscript{(94)} বা \VUgras{বাছুরের} \textsuperscript{(92)} --- \VUgras{রান,
চাপ, টুকরো} ও \VUgras{পাঁজরের} আকারে। \VUgras{কসাইয়ের ছোকরা}
\textsuperscript{(93)} এই সমস্ত মাংস কাটে ও \VUgras{মজ্জার হাড়}
\textsuperscript{(96)} আলাদা রাখে। কসাইয়ের দরজায় এক \VUgras{ঝিনুক-বিক্রেতা}
\textsuperscript{(97)} দাঁড়িয়ে, যে \VUgras{ঝিনুক} \textsuperscript{(98)} বেচে।
চাইলে সে সেগুলি খুলেও দেয়।

%%K t15-16-1 p
16. --- এই সব ব্যবসায়ী লাইসেন্স বা জায়গার শুল্ক দেন; তাই \VUgras{সিপাহি}
\textsuperscript{(100)} সেই বেচারা \VUgras{ঘুরে তরকারি বেচা মেয়েদের}
\textsuperscript{(99)} নির্দয়ভাবে তাড়া করে, যারা নিজেদের ঠেলায় \VUgras{গাজর,
মুলো, শালগম, পেঁয়াজকলি} বা \VUgras{শতমূলীর আঁটি, টাটকা মটর, পালং, চুকা,
জলকুম্ভি} ও \VUgras{ধনেপাতা} বয়ে তাঁদের সঙ্গে পাল্লা দেয়।

%%K t15-tit-8 sub
\VUsaut{4.0mm}
\VUpk{10.2pt}{40}{সিপাহি থেকে সাবধান!}
\VUsaut{3.0mm}

%%K t15-17-1 p
17. --- \VUgras{রসুন} \textsuperscript{(101)} ও \VUgras{পেঁয়াজ} বেচা ছেলেটি বেশ
জানে যে তারও বাজারের কাছে নিজের মাল বেচার অনুমতি নেই। সে ছুটে পালায়, এই ঝুঁকি
নিয়ে যে \VUgras{গলদা চিংড়ি} ও \VUgras{লবস্টার} \VUgras{বিক্রেতার}
\textsuperscript{(102)} \VUgras{কাঁকড়া}, \VUgras{চিংড়ি-মাছ} ও \VUgras{ছোট
চিংড়ির} ঝুড়ি উলটে দেয়।

%%K t15-18-1 p
18. --- সবাই তাড়াহুড়োয় আর বড় হইচই করছে; কিন্তু এতে ঘুরন্ত
\VUgras{কাচমিস্ত্রির} \textsuperscript{(103)} হাঁক স্পষ্ট শোনা থামে না, সেই
\VUgras{কয়লাওয়ালারও} \textsuperscript{(104)} নয়, যে এইমাত্র একটি
\VUgras{কয়লার বস্তা} \textsuperscript{(105)} পৌঁছে দিতে নিজের গাড়ি কিছুক্ষণের
জন্য ছেড়েছে।
\VUsaut{6.0mm}
\VUfilet{22mm}
